\section{Uvod}
\label{sec:uvod}

Razvoj sustava automatizirane vožnje (engl.\ \textit{Automated Driving Systems}, ADS) predstavlja jednu od najznačajnijih tehnoloških transformacija u automobilskoj industriji. Prema klasifikaciji SAE J3016 \cite{saej3016}, razine automatizacije kreću se od razine~0 (potpuno ručno upravljanje) do razine~5 (puna automatizacija bez ikakve potrebe za ljudskom intervencijom). Sa svakom višom razinom automatizacije, sustav preuzima sve veću odgovornost za dinamičke zadatke vožnje, čime se istodobno povećava i složenost osiguravanja sigurnosti takvih sustava.

U kontekstu sigurnosti automobilskih električno-elektroničkih (E/E) sustava, dvije norme Međunarodne organizacije za normizaciju (ISO) imaju ključnu ulogu:
\begin{itemize}
  \item \textbf{ISO~26262} \cite{iso26262} — norma za \textit{funkcijsku sigurnost} (engl.\ \textit{Functional Safety}, FuSa), koja se bavi sprečavanjem rizika uzrokovanih kvarovima i sustavnim pogreškama u E/E komponentama;
  \item \textbf{ISO~21448} \cite{iso21448} — norma za \textit{sigurnost namijenjene funkcionalnosti} (engl.\ \textit{Safety of the Intended Functionality}, SOTIF), koja adresira opasnosti koje nastaju zbog ograničenja performansi sustava, čak i kada sustav radi točno onako kako je projektiran.
\end{itemize}

Dok je ISO~26262 razvijena kao proširenje generičke norme IEC~61508 specifično za cestovna vozila, ISO~21448 je novija norma nastala iz prepoznavanja da funkcijska sigurnost sama po sebi nije dovoljna za sustave koji koriste napredne senzore, algoritme strojnog učenja i kompleksne strategije donošenja odluka. Upravo ta komplementarnost dviju normi čini ih nezaobilaznim temeljem sigurnosnog okvira za automatizirane vozne sustave.

Cilj ovog seminara je sustavno prikazati strukturu, ključne koncepte i metodologije obiju normi, usporediti njihove pristupe upravljanju rizikom te raspraviti izazove njihove integracije u cjeloviti razvojni proces sustava automatizirane vožnje.

Rad je organiziran na sljedeći način: u drugom poglavlju detaljno se obrađuje norma ISO~26262, uključujući V-model razvoja, HARA metodologiju i ASIL klasifikaciju. Treće poglavlje posvećeno je normi ISO~21448 i konceptu SOTIF-a. Četvrto poglavlje donosi usporednu analizu i raspravu o integraciji dviju normi, dok peto poglavlje zaključuje rad sintezom nalaza i pogledom na buduće smjerove razvoja normativnog okvira.
